%
%  Example solution set
%
\documentclass[11pt]{article}
% 
\usepackage{amsmath}
\usepackage{amssymb}
\usepackage{graphicx}
\usepackage{float}
\usepackage{calc}
\usepackage[margin=1.in]{geometry}
\usepackage{bm}
\usepackage[usenames]{color}

\usepackage{fancyvrb}% fancy Verbatim
\usepackage{media9} % to attach viedo file in latex
\usepackage{listings}
\lstset{
	basicstyle=\footnotesize\ttfamily,
	columns=flexible,
	breaklines=true,
	tabsize=2,
	commentstyle=\color{red},
	keywordstyle=\color{blue}\bfseries,
	keepspaces=true
}

\lstset{language=C++,
	basicstyle=\footnotesize\ttfamily,
	keywordstyle=\color{blue}\footnotesize\ttfamily,
	stringstyle=\color{green}\footnotesize\ttfamily,
	commentstyle=\color{red}\footnotesize\ttfamily,
	morecomment=[l][\color{magenta}]{\#},
	columns=flexible,
	breaklines=true,
	keepspaces=true        
}
\def\ab{{\bm a}} \def\cb{{\bm c}} \def\db{{\bm d}} \def\eb{{\bm e}}
\def\bb{{\bm b}} \def\yb{{\bm y}}
\def\ub{{\bm u}}  \def\xb{{\bm x}} \def\vb{{\bm v}} \def\wb{{\bm w}}
\def\omegab{{\bm \omega}} \def\rb{{\bm r}} \def\ib{{\bm i}} \def\jb{{\bm j}}
\def\lb{{\bm l}} \def\kb{{\bm k}} \def\Ab{{\bm A}} \def\fb{{\bm f}} \def\Ub{{\bm U}}
\def\Fb{{\bm F}} \def\nb{{\bm n}} \def\Db{{\bm D}} \def\eb{{\bm e}} 
\def\gb{{\bm g}}  \def\Gb{{\bm G}} \def\hb{{\bm h}} \def\Yb{{\bm Y}} \def\Rb{{\bm R}} 
\def\Tb{{\bm T}} \def\zb{{\bm z}} \def\Zb{{\bm Z}} \def\Qb{{\bm Q}} \def\Cb{{\bm C}} 
\def\Bb{{\bm B}} \def\Rbb{{\mathbb R}} \def\Cbb{{\mathbb C}} \def\Lb{{\bm L}} \def\Mb{{\bm M}}
\def\Hb{{\bm H}} \def\Wb{{\bm W}} \def\Vb{{\bm V}}
\newcommand{\Dp}{D_+}
\newcommand{\Dm}{D_-}
\newcommand{\Dz}{D_0}
\newcommand{\dx}{\Delta x}
\newcommand{\dy}{\Delta y}
\newcommand{\dt}{\Delta t}
\newcommand{\Nx}{N_x} 
\newcommand{\norm}[1]{\left\Vert{#1}\right\Vert}
\begin{document}
	
	
\begin{center}
\Large{MATH 6890: PARALLEL METHODS FOR SCIENTIFIC COMPUTING \\
Fall 2023}\\
\textbf{Instructor:} Prof. Henshaw
\end{center}
\begin{center}
 Ngan Le\\
\textbf {FINAL PROJECT \\
PARALLEL SOLVER FOR ANISOTROPIC WAVE EQUATION} 
\end{center}

	
	\medskip
	% -------------------------
	
	\section{Project description}
\subsection{Overview}
The aim of this project is to develop parallel solvers for the anisotropic wave equation in two-dimensional space with Dirichlet boundary conditions on a rectangular domain $\Omega = (a_x,b_x) \times (a_y,b_y)$
\begin{subequations}
\label{eq:2DAnisotropicWaveEquation}
\begin{align}
& \rho(x,y)w_{tt}=\nabla\cdot\left(\bm{m}(x,y)\cdot\nabla w\right) +f(x,y,t), && (x,y)\in \Omega,\; 0 \leq t \leq t_{\text{final}}, \label{eq:2DAnisotropicWaveEquationA} \\
& w(a_x,y,t) = g_{ax}(y,t), \quad w(b_x,y,t) = g_{bx}(y,t), && y \in [a_y,b_y],\; 0 \leq t \leq t_{\text{final}}, \label{eq:2DAnisotropicWaveEquationB} \\
& w(x,a_y,t) = g_{ay}(x,t), \quad w(x,b_y,t) = g_{by}(x,t), && x \in [a_x,b_x],\; 0 \leq t \leq t_{\text{final}}, \label{eq:2DAnisotropicWaveEquationC} \\
& w(x,y,0) = v_0(x,y), \quad w_t(x,y,0) = v_1(x,y), && (x,y) \in \Omega. \label{eq:2DAnisotropicWaveEquationD}
\end{align}
\end{subequations}
Here $\rho(x,y) > 0$ is the density and $\bm{m}(x,y)$ is the stiffness matrix (material stiffness tensor),
\begin{align*}
& \bm{m}(x,y)=\begin{bmatrix}
	m_{11}(x,y) & m_{12}(x,y)\\
	m_{21}(x,y) & m_{22}(x,y)
\end{bmatrix}.
\end{align*}
The matrix $\bm{m}$ is symmetric ($m_{12}=m_{21}$) and positive definite for all $(x,y)\in \Omega$. Note that $\rho$ and $\bm{m}$ are piecewise continuous on $\Omega$. The derivatives in~\eqref{eq:2DAnisotropicWaveEquationA} can be written explicitly as
\begin{align}
\nabla\cdot\left(\bm{m}\cdot\nabla w\right) = \partial_x (m_{11}w_x + m_{12}w_y) + \partial_y (m_{21}w_x + m_{22}w_y). \label{eq:2DAnisotropicOperator}
\end{align}

\subsection{Numerical scheme}
The grid is $x_i = a_x + i\Delta x$, $y_j = a_y + j\Delta y$, $t^n = n\Delta t$, with $\Delta x = (b_x-a_x)/N_x$, $\Delta y = (b_y-a_y)/N_y$ and $\Delta t = t_{\text{final}}/N_t$. The diagonal coefficients are averaged to the half grid points, and the off-diagonal coefficients are evaluated at the grid points,
\begin{align*}
& M^{11}_{i\pm1/2,j} = \dfrac{m_{11}(x_{i},y_j)+m_{11}(x_{i\pm1},y_j)}{2}, \quad M^{22}_{i,j\pm1/2} = \dfrac{m_{22}(x_{i},y_j)+m_{22}(x_{i},y_{j\pm1})}{2}, \\
& M^{12}_{ij} = m_{12}(x_i,y_j), \quad M^{21}_{ij} = m_{21}(x_i,y_j),
\end{align*}
and we write $f_{ij}^n=f(x_i,y_j,t^n)$, $\rho_{ij}=\rho(x_i,y_j)$. The second-order accurate finite difference approximation of the operator~\eqref{eq:2DAnisotropicOperator} is
\begin{align}
L_h W_{ij} ={}& D_{+x}\left[M^{11}_{i-1/2,j}D_{-x}W_{ij}\right]+D_{+y}\left[M^{22}_{i,j-1/2}D_{-y}W_{ij}\right] \notag\\
& + D_{0x}\left[M^{12}_{ij}D_{0y}W_{ij}\right]+D_{0y}\left[M^{21}_{ij}D_{0x}W_{ij}\right]. \label{eq:2DAnisotropicLh}
\end{align}
\textcolor{red}{\textbf{Note (current C++ code, wrong):} the OpenMP and MPI codes (fixed in the serial code) average the off-diagonal coefficients to the half grid points but keep the differences at $i\pm1$ and $j\pm1$,
\begin{align*}
& \dfrac{1}{2\Delta x}\left[M^{12}_{i+1/2,j}D_{0y}W_{i+1,j} - M^{12}_{i-1/2,j}D_{0y}W_{i-1,j}\right] \\
& + \dfrac{1}{2\Delta y}\left[M^{21}_{i,j+1/2}D_{0x}W_{i,j+1} - M^{21}_{i,j-1/2}D_{0x}W_{i,j-1}\right], \\
& M^{12}_{i\pm1/2,j} = \dfrac{m_{12}(x_{i\pm1},y_j)+m_{12}(x_i,y_j)}{2}, \quad M^{21}_{i,j\pm1/2} = \dfrac{m_{21}(x_i,y_{j\pm1})+m_{21}(x_i,y_j)}{2}.
\end{align*}
This is inconsistent when $m_{12}$ depends on $x$ (or $m_{21}$ on $y$): for $m_{12}=x$ and $D_{0y}W=1$ it gives $(M^{12}_{i+1/2,j}-M^{12}_{i-1/2,j})/(2\Delta x)=1/2$ instead of $\partial_x m_{12}=1$. To be replaced by the pointwise form in~\eqref{eq:2DAnisotropicLh}.}

The explicit scheme for~\eqref{eq:2DAnisotropicWaveEquation} is given by
\begin{subequations}
\label{eq:explicitScheme}
\begin{enumerate}
\item [$\bullet$] \textcolor{blue}{Difference equation:}
\begin{align}
& \rho_{ij}\, D_{+t}D_{-t}W_{ij}^n = L_h W_{ij}^n + f_{ij}^n, \label{eq:explicitSchemeA}\\
& \quad i=1,2,\dots,N_x-1,\quad j=1,2,\dots, N_y-1,\quad n=1,2,\dots, N_t-1. \notag
\end{align}

\item [$\bullet$] \textcolor{blue}{Initial conditions:}
\begin{align}
& W_{ij}^0=V_{0,ij}, && i=1,\dots,N_x-1,\; j=1,\dots, N_y-1, \label{eq:explicitSchemeB}\\
& W_{ij}^1=V_{0,ij}+\Delta t\, V_{1,ij} +\dfrac{\Delta t^2}{2\rho_{ij}}\left[L_h V_{0,ij} + f_{ij}^0\right], && i=1,\dots,N_x-1,\; j=1,\dots, N_y-1, \label{eq:explicitSchemeC}
\end{align}
where $V_{0,ij} = v_0(x_i,y_j)$ and $V_{1,ij} = v_1(x_i, y_j)$.

\item [$\bullet$] \textcolor{blue}{Boundary conditions:}
\begin{align}
&W_{0j}^n = g_{ax}(y_j,t^n), \quad W_{N_x j}^n = g_{bx}(y_j,t^n), && j=0,1,\dots, N_y,\; n=0,1,\dots, N_t, \label{eq:explicitSchemeD}\\
&W_{i0}^n = g_{ay}(x_i,t^n), \quad W_{i N_y}^n = g_{by}(x_i,t^n), && i=0,1,\dots, N_x,\; n=0,1,\dots, N_t. \label{eq:explicitSchemeE}
\end{align}
\end{enumerate}
\end{subequations}
The approximation for $W^1_{ij}$ is obtained from the Taylor series expansion of $w(x,y,\Delta t)$, together with the PDE and the initial conditions:
\begin{align*}
w(x,y,\Delta t) &= w(x,y,0)+\Delta t\, w_t(x,y,0) + \dfrac{\Delta t^2}{2} w_{tt}(x,y,0)+\mathcal{O}(\Delta t^3)\\
	&=v_0(x,y) + \Delta t\, v_1(x,y) + \dfrac{\Delta t^2}{2\rho(x,y)} \left[\nabla\cdot\left(\bm{m}(x,y)\cdot\nabla v_0\right) +f(x,y,0)\right]+\mathcal{O}(\Delta t^3).
\end{align*}
Thus, substituting $V_{0,ij}$ and $V_{1,ij}$ gives the approximation in~\eqref{eq:explicitSchemeC}.

The time step restriction is
\begin{align}
	\Delta t \leq \mbox{cfl} \cdot \min_{i,j}\dfrac{\sqrt{\rho_{ij}}}{\sqrt{\dfrac{m_{11}(x_i,y_j)}{\Delta x^2} + \dfrac{\left|m_{12}(x_i,y_j)+m_{21}(x_i,y_j)\right|}{4\Delta x \Delta y} + \dfrac{m_{22}(x_i,y_j)}{\Delta y^2}}}, \quad \mbox{cfl} < 1.
\end{align}
The number of time steps is then $N_t = \lceil t_{\text{final}}/\Delta t \rceil$, and $\Delta t$ is reset to $t_{\text{final}}/N_t$ so that the final time is reached exactly.



\section{Numerical results}
\subsection{Test cases}
All tests are on $\Omega=(-1,1)\times(-1,1)$ with $t_{\text{final}}=0.5$, $\mbox{cfl}=0.9$ and $N_x=N_y=20,\,40,\,80,\,160$. Two sets of coefficients are used:
\begin{itemize}
\item \textcolor{blue}{Constant coefficients (isotropic material):} $\rho=1$, $\bm{m}=I$.
\item \textcolor{blue}{Variable coefficients:}
\begin{align*}
\rho(x,y) = 0.1\left(\sin(x+y)+1.1\right), \quad
\bm{m}(x,y)=\begin{bmatrix}
	0.3(x+y)+2 & 0.5+0.2(x+y)\\
	0.5+0.2(x+y) & 0.6(2x+y)+3
\end{bmatrix}.
\end{align*}
The matrix $\bm{m}$ is symmetric ($m_{12}=m_{21}$) and positive definite on $\Omega$: $m_{11}\geq 1.4$ and $m_{11}m_{22}-m_{12}^2\geq 1.68-0.81>0$.
\end{itemize}
The exact solutions are
\begin{itemize}
\item \textcolor{blue}{poly} (manufactured solution, $f=\rho w_{tt}-\nabla\cdot(\bm{m}\cdot\nabla w)$):
$w=(b_0+b_1x+b_2x^2)(c_0+c_1y+c_2y^2)(a_0+a_1t+a_2t^2)$, with $(a_0,a_1,a_2)=(0.5,0.7,0.9)$, $(b_0,b_1,b_2)=(0.5,0.8,0.6)$, $(c_0,c_1,c_2)=(0.3,0.25,0.2)$.
\item \textcolor{blue}{trig}: $w=3\cos(s)$,
\quad \textcolor{blue}{pulse}: $w=4e^{-3s^2}$,
\quad \textcolor{blue}{trigPulse}: $w=2e^{-2.5s^2}\cos(2s)$,
\end{itemize}
where $s=k_xx+k_yy-\omega t$, $k_x=k_y=1$, $f=0$, and $\omega$ satisfies the dispersion relation $\rho\,\omega^2=m_{11}k_x^2+(m_{12}+m_{21})k_xk_y+m_{22}k_y^2$. The plane waves (trig, pulse, trigPulse) are exact solutions only for constant coefficients. The five test cases are
\begin{table}[H]\footnotesize
\begin{center}
\begin{tabular}{|c|c|c|l|} \hline
Case & Coefficients & Solution & Expected \\ \hline
1 & constant & poly      & exact (round-off) \\
2 & constant & trig      & second order (ratio $\approx 4$) \\
3 & constant & pulse     & second order \\
4 & constant & trigPulse & second order \\
5 & variable & poly      & second order (not exact: $m_{12}$ is linear in $x$, so $m_{12}w_y$ is cubic in $x$) \\ \hline
\end{tabular}
\end{center}
\end{table}
The MATLAB baseline is run with \texttt{matlab/runBaselineMATLAB.m} and the serial baseline with \texttt{serial/runBaselineSerial.sh}; the results are saved in the folder \texttt{reference}. A run passes if the number of time steps is the same and
\begin{align*}
\max_{N_x}\left|\mbox{maxErr}-\mbox{maxErr}_{\text{ref}}\right|<10^{-12}.
\end{align*}

\subsection{Serial version}
The results below are the output of the serial baseline script. The four lines of each case are $N_x=N_y=20,\,40,\,80,\,160$ at the final time $t_{\text{final}}=0.5$, and \texttt{ratio} is the ratio of the errors on two successive grids.

$\bullet$ Case 1: constant coefficients, poly (exact up to round-off)
\begin{Verbatim}[fontsize=\footnotesize,formatcom=\color{blue},frame=single,framerule=1pt,numbers=left]
poly: Nt=  8 dt=6.25e-02 maxNorm=1.53e+00 maxRelErr=2.00e-16 cpuTime= 9.00e-05(s)
poly: Nt= 16 dt=3.12e-02 maxNorm=1.53e+00 maxRelErr=2.68e-16 cpuTime= 1.79e-04(s) ratio=0.75
poly: Nt= 32 dt=1.56e-02 maxNorm=1.53e+00 maxRelErr=4.46e-16 cpuTime= 1.79e-03(s) ratio=0.60
poly: Nt= 63 dt=7.94e-03 maxNorm=1.53e+00 maxRelErr=7.19e-16 cpuTime= 9.89e-03(s) ratio=0.62
\end{Verbatim}

$\bullet$ Case 2: constant coefficients, trig
\begin{Verbatim}[fontsize=\footnotesize,formatcom=\color{blue},frame=single,framerule=1pt,numbers=left]
trig: Nt=  8 dt=6.25e-02 maxNorm=3.00e+00 maxRelErr=6.74e-04 cpuTime= 2.30e-05(s)
trig: Nt= 16 dt=3.12e-02 maxNorm=3.00e+00 maxRelErr=1.75e-04 cpuTime= 1.01e-04(s) ratio=3.86
trig: Nt= 32 dt=1.56e-02 maxNorm=3.00e+00 maxRelErr=4.47e-05 cpuTime= 6.25e-04(s) ratio=3.91
trig: Nt= 63 dt=7.94e-03 maxNorm=3.00e+00 maxRelErr=1.16e-05 cpuTime= 4.22e-03(s) ratio=3.84
\end{Verbatim}

$\bullet$ Case 3: constant coefficients, pulse
\begin{Verbatim}[fontsize=\footnotesize,formatcom=\color{blue},frame=single,framerule=1pt,numbers=left]
pulse: Nt=  8 dt=6.25e-02 maxNorm=4.02e+00 maxRelErr=5.53e-03 cpuTime= 4.00e-05(s)
pulse: Nt= 16 dt=3.12e-02 maxNorm=4.00e+00 maxRelErr=1.41e-03 cpuTime= 1.13e-04(s) ratio=3.93
pulse: Nt= 32 dt=1.56e-02 maxNorm=4.00e+00 maxRelErr=3.50e-04 cpuTime= 7.24e-04(s) ratio=4.04
pulse: Nt= 63 dt=7.94e-03 maxNorm=4.00e+00 maxRelErr=8.87e-05 cpuTime= 4.08e-03(s) ratio=3.95
\end{Verbatim}

$\bullet$ Case 4: constant coefficients, trigPulse
\begin{Verbatim}[fontsize=\footnotesize,formatcom=\color{blue},frame=single,framerule=1pt,numbers=left]
trigPulse: Nt=  8 dt=6.25e-02 maxNorm=2.02e+00 maxRelErr=8.88e-03 cpuTime= 4.10e-05(s)
trigPulse: Nt= 16 dt=3.12e-02 maxNorm=2.00e+00 maxRelErr=2.20e-03 cpuTime= 1.15e-04(s) ratio=4.06
trigPulse: Nt= 32 dt=1.56e-02 maxNorm=2.00e+00 maxRelErr=5.50e-04 cpuTime= 7.83e-04(s) ratio=4.01
trigPulse: Nt= 63 dt=7.94e-03 maxNorm=2.00e+00 maxRelErr=1.39e-04 cpuTime= 4.35e-03(s) ratio=3.97
\end{Verbatim}

$\bullet$ Case 5: variable coefficients, poly
\begin{Verbatim}[fontsize=\footnotesize,formatcom=\color{blue},frame=single,framerule=1pt,numbers=left]
poly: Nt=102 dt=4.90e-03 maxNorm=1.53e+00 maxRelErr=4.18e-05 cpuTime= 5.56e-04(s)
poly: Nt=204 dt=2.45e-03 maxNorm=1.53e+00 maxRelErr=1.04e-05 cpuTime= 3.72e-03(s) ratio=4.00
poly: Nt=409 dt=1.22e-03 maxNorm=1.53e+00 maxRelErr=2.61e-06 cpuTime= 4.76e-02(s) ratio=4.00
poly: Nt=817 dt=6.12e-04 maxNorm=1.53e+00 maxRelErr=6.53e-07 cpuTime= 3.73e-01(s) ratio=4.00
\end{Verbatim}

The serial code agrees with the MATLAB code (\texttt{matlab/anisotropicWave\_2D.m}) up to round-off:
\begin{table}[H]\footnotesize
\begin{center}
\begin{tabular}{|c|c|c|c|} \hline
Case & Coefficients & Solution & $\max|\mbox{maxErr}_{\text{C++}}-\mbox{maxErr}_{\text{MATLAB}}|$ \\ \hline
1 & constant & poly      & $4.28\cdot 10^{-16}$ \\
2 & constant & trig      & $2.44\cdot 10^{-15}$ \\
3 & constant & pulse     & $2.22\cdot 10^{-15}$ \\
4 & constant & trigPulse & $1.53\cdot 10^{-15}$ \\
5 & variable & poly      & $2.10\cdot 10^{-15}$ \\ \hline
\end{tabular}
\end{center}
\end{table}
\subsection { Parallel version}
\textcolor{red}{Results from the old code (before the fixes); to be updated.}\\
Here we are testing for polynomial manufactured solution with constant coefficients using 4 threads/processes. We observe that the errors are the same as those in the serial code.\\
	$\bullet$ OpenMP code
%			\begin{table}[H]\footnotesize 
%					\begin{center}
%						\begin{tabular}{|c|c|c|c|c|} \hline 
%						%	\multicolumn{4}{|c|}{heat1d CUDA version}    \\  \hline
%							$\mbox{nx}$   & $\mbox{tFinal}$        & $\mbox{Nt}$   & $\mbox{maxNorm}$  & $\mbox{maxRelErr}$	\\  \hline 
%							$ 20 $     & $\mbox{0.5}$          & $ 21 $       & $1.53e+00$   & $3.62e-16$   \\ 
%							$ 40 $     & $\mbox{0.5}$          & $ 41 $       & $1.53e+00$   & $5.80e-16$  \\ 
%							$ 80 $    & $\mbox{0.5}$          & $ 81 $       & $1.53e+00$   &  $8.70e-16$  \\ 
%							$ 160 $   & $\mbox{0.5}$         & $ 161 $       & $1.53e+00$   &   $1.59e-15$  \\ 
%							\hline
%						\end{tabular}
%						\caption{
%							Results for polynomial manufactured solution with 1 thread
%						}
%						\label{tab:openMPwave2d}
%					\end{center}
%				\end{table}
\begin{Verbatim}[fontsize=\footnotesize,formatcom=\color{blue},frame=single,framerule=1pt,numbers=left]
./wave2dOpenMP -nx=20 -tFinal=0.5 -numThreads=4 -debug=0
nx=20 ny=20 dt=2.38e-02 numSteps=21 tFinal=  0.50
CPU: numSteps=21 nx=20 maxNorm=1.53e+00 maxRelErr=3.62e-16 cpuTime= 7.62e-04(s)
	
./wave2dOpenMP -nx=40 -tFinal=0.5 -numThreads=4 -debug=0
nx=40 ny=40 dt=1.22e-02 numSteps=41 tFinal=  0.50
CPU: numSteps=41 nx=40 maxNorm=1.53e+00 maxRelErr=5.80e-16 cpuTime= 5.03e-03(s)
	
./wave2dOpenMP -nx=80 -tFinal=0.5 -numThreads=4 -debug=0
nx=80 ny=80 dt=6.17e-03 numSteps=81 tFinal=  0.50
CPU: numSteps=81 nx=80 maxNorm=1.53e+00 maxRelErr=8.70e-16 cpuTime= 3.90e-02(s)
\end{Verbatim}
\newpage
	$\bullet$ MPI code
\begin{Verbatim}[fontsize=\footnotesize,formatcom=\color{blue},frame=single,framerule=1pt,numbers=left]
mpirun -np=4 ./wave2dMPI -tFinal=.5 -nx=20 -debug=0
nx=20 ny=20 dt=2.38e-02 numSteps=21 tFinal=  0.50
np=4 myRank=0 Gathered solution: maxNormG=1.53e+00 maxRelErrG=3.62e-16 cpuTimeStepG= 2.26e-04(s)

mpirun -np=4 ./wave2dMPI -tFinal=.5 -nx=40 -debug=0
nx=40 ny=40 dt=1.22e-02 numSteps=41 tFinal=  0.50
np=4 myRank=0 Gathered solution: maxNormG=1.53e+00 maxRelErrG=5.80e-16 cpuTimeStepG= 6.04e-04(s)

mpirun -np=4 ./wave2dMPI -tFinal=.5 -nx=80 -debug=0
nx=80 ny=80 dt=6.17e-03 numSteps=81 tFinal=  0.50
np=4 myRank=0 Gathered solution: maxNormG=1.53e+00 maxRelErrG=8.70e-16 cpuTimeStepG= 7.57e-03(s)
\end{Verbatim}


\section{Parallel speedup.}
\textcolor{red}{Results from the old code (before the fixes); to be updated.}\\
For testing speedup, we run the code on cg66 using $nx = ny = 2048$, $\mbox{tFinal=0.1}$. The serial code took 32.9 seconds to run. Here are the results for openMP and MPI versions.
\begin{table}[H]\footnotesize 
\begin{center}
\begin{tabular}{|c|c|c|c|} \hline 
\multicolumn{4}{|c|}{openMP version}    \\  \hline
$\mbox{Threads}$   & $\mbox{cpu (s)}$        & $\mbox{Speedup}$   & $\mbox{Par-Eff}$   	\\  \hline 
$ 1 $     & $34.2$          & $ 1 $               & $1$       \\ 
$ 2$     & $\mbox{17.6}$          & $ 1.94 $       & $0.74$    \\ 
$ 4$    & $\mbox{9.71}$          & $ 3.52 $       & $0.88$   \\ 
$ 8 $   & $\mbox{4.76}$         & $ 7.18 $       & $0.89$   \\ 
$ 16 $   & $\mbox{2.48}$         & $ 13.79 $       & $0.86$   \\
$ 32 $   & $\mbox{1.27}$         & $ 26.92 $       & $0.84$   \\
\hline
\end{tabular}
\caption{
Speedup for openMP code.
}
\label{tab:openMP code}
\end{center}
\end{table}

\begin{table}[H]\footnotesize 
\begin{center}
\begin{tabular}{|c|c|c|c|} \hline 
\multicolumn{4}{|c|}{MPI version}    \\  \hline
$\mbox{np}$   & $\mbox{cpu (s)}$        & $\mbox{Speedup}$   & $\mbox{Par-Eff}$   	\\  \hline 
$ 1 $     & $32.5$          & $ 1 $       & $1$       \\ 
$ 2$     & $\mbox{16.6}$          & $ 1.95 $       & $0.97$    \\ 
$ 4$    & $\mbox{9.09}$          & $ 3.34 $       & $0.83$   \\ 
$ 8 $   & $\mbox{4.81}$         & $ 6.75 $       & $0.84$   \\ 
$ 16 $   & $\mbox{2.52}$         & $ 12.89 $       & $0.80$   \\
$ 32 $   & $\mbox{2.17}$         & $ 14.97 $       & $0.46$   \\
\hline
\end{tabular}
\caption{
Speedup for MPI code.
}
\label{tab:MPI code}
\end{center}
\end{table}
\clearpage




	
	
	
% \lstinputlisting[language=C++, numbers=left, stepnumber=1, firstline=1,caption={wave2dSerial.C},label=code:serial,frame=single]{"D:/Documents/Graduate/PHD/RPI/Fall 2023/MATH-6890 Parallel Methods/Final Project/serialCode/wave2dSerial.C"}
% \lstinputlisting[language=C++, numbers=left, stepnumber=1, firstline=1,caption={wave2dOpenMP.C},label=code:openMP,frame=single]{"D:/Documents/Graduate/PHD/RPI/Fall 2023/MATH-6890 Parallel Methods/Final_Project/openMP Code/wave2dOpenMP.C"}
% \lstinputlisting[language=C++, numbers=left, stepnumber=1, firstline=1,caption={wave2dMPI.C},label=code:MPI,frame=single]{"D:/Documents/Graduate/PHD/RPI/Fall 2023/MATH-6890 Parallel Methods/Final Project/MPI Code/wave2dMPI.C"}
%	\lstinputlisting[language=C++, numbers=left, stepnumber=1, firstline=1,caption={heat2d.cu},label=code:heat2d,frame=single]{"D:/Documents/Graduate/PHD/RPI/Fall 2023/MATH-6890 Parallel Methods/CUDA code/heat2d.C"}
%	\lstinputlisting[language=C++, numbers=left, stepnumber=1, firstline=1,caption={heatADI.C},label=code:Makefile,frame=single]{code/heatADI.C}
	%\lstinputlisting[language=C++, numbers=left, stepnumber=1, firstline=1,caption={heat2dUpdate.f90},label=code:heat2dUpdate,frame=single]{code/heat2dUpdate.f90}
	%\lstinputlisting[language=C++, numbers=left, stepnumber=1, firstline=1,caption={heat1d.C},label=code:heat1dParallelVersion,frame=single]{code/heat1d.C}
	
	%\lstinputlisting[language=Matlab, numbers=left, stepnumber=1, firstline=1,caption={writeMatlabArray.h},label=code:writeMatlabArray,frame=single]{code/writeMatlabArray.h}
	
%	\lstinputlisting[language=C++, numbers=left, stepnumber=1, firstline=1,caption={poisson.C},label=code:poissonParallelVersion,frame=single]{code/poisson.C}
\end{document}
% =====================================================
% =====================================================
% =====================================================




